\subsection{The remaining return sectors and the exact outstanding attachment}
\begin{statement}{P3K.75}{Unresolved; exact remaining obligations}
The complete scalar-fan algorithm of P3K.73 and selected-cusp
source-fan algorithm of P3K.71 settle their stated marked loops.
The remaining source-zero return is specified by the actual family
\[
 Y^2=X^3+(1-2v)X^2+TX+v^2,\quad G=(0,v),\quad
 q=T-(v^2-2v),\quad \lambda=v^{km}\lambda_0,
\]
with its retained cap data
\[
 (a,b,m)=(1,2,4p)\quad\hbox{or}\quad(2,1,4p-1),
 \qquad p,k\ge1,\quad n=\pm1,
\]
including the marked Cartier line, finite character and complete
annular cochain. On the initial and reflected fans the proved
constraints are precisely
\[
 (Z_{\mathcal T},Z_{J\mathcal T})=
 (kA_{\rm scalar}+14e,a+kA_{\rm scalar}+14e),\qquad
 e\in\{0,1\},
\]
where \(A_{\rm scalar}=0\) on the first branch and
\(A_{\rm scalar}=3p-1\) on the second. Determining \(e\)
from the actual oriented global attachment remains an obligation.
The smallest cases are \((0,1)\) or \((14,15)\), and
\((2,4)\) or \((16,18)\), respectively.

K77 now computes the component-supported OS49 return on every
lattice-basis fan, including the actual shorter endpoint and
the retained fixed-source marked cap-data iterates. Its exact
image on either specified parameter loop is
\(\gcd(k,28)\mathbb Z/28\mathbb Z\).
Other remaining parts of item 3.4 are the added ordinary-clutch
sectors, the remaining marked source-moduli loops, and unmarked
endpoint sectors. No value is assigned here to those sectors.
\end{statement}
\begin{proof}[Recorded reductions, attempted completions and the next exact maps]
\par\medskip\noindent\textbf{The source-zero integer that remains.}\par\nopagebreak[4]
The physical/dual line map and both Hopf-gluing orientations are
proved in K68. K69 carries them through the actual global
parameter groupoid and the full source-infinity cap comparison.
They give the paired sum. K60's retained single native diagonal
pair \((1-a,a)\), with the literal single-to-double pullback
product, gives the difference \(Z_{J\mathcal T}-Z_{\mathcal T}=a\).
Solving these two equations in the stated group leaves exactly
the kernel \(\{0,14\}\) of multiplication by two. These actual
maps prove the displayed two possibilities; choosing either
element requires further information.

The completed selected-cusp and vertical-source calculations
do not supply that further integer: their loops and endpoints
are the explicitly different paths recorded in the global
groupoid. Their composition yields the paired relation just
used. The precise groupoid arrows, rather than an identification
of individual loops, are retained in K58--60 and K69.

\par\medskip\noindent\textbf{The attempted source completion and its actual obstruction.}\par\nopagebreak[4]
The marked small-resolution source charts and both Cartier
spinor maps have been constructed in the source-zero completion
notes. On the unresolved central contraction the cyclic
generator fixes the cone origin; every power fixes it. Thus
the literal origin stabilizer is infinite and that proposed
central quotient does not meet the proper-discontinuity
condition used for the punctured construction. Deleting or
resolving this point changes the required gluing problem.
No index for a smooth isolated quotient has been assigned
to that infinite-stabilizer point.

For the earlier specified fractional-period root quotient
of SZ.27, the root map has local base equation \(q=z^m\);
its central
divisor therefore has multiplicity \(m\), rather than the
specified simple selected zero. Moving that branching to a
regular value inserts its actual additional cap. Those maps
are retained with these exact changes of the requested data.
The subsequently constructed global noninvariant norm sections
have their proved source identities, but their full line and
root action has not been constructed. The multiplicity
calculation just given applies to the earlier SZ.27 quotient;
no multiplicity is assigned to the new source-section route.
Its full source-line-cap comparison is still required.

\par\medskip\noindent\textbf{The real trace and its extra end.}\par\nopagebreak[4]
The current actual normalization trace is
\(F=\operatorname{Bl}_{(-1,0)}(\mathbb P^1_s\times D_v)\).
Its two blowup-chart maps are written in
\path{notes/part3_source_zero_real_trace_working.md}; their
Jacobians are units on the exceptional curve. The actual
source line restricts to \(\mathcal O(-aE)\), with
\(E^2=-1\). The two transported spinor determinants are
\((1-2a)E\) and \(-(1+2a)E\), giving local intersection
indices \(a(a-1)/2\) and \(a(a+1)/2\). Their difference
agrees with the proved native difference \(a\).

The exact Cox coordinates also give
\[
 k=s_1-s_0,\qquad l=t_0+t_1,\qquad
 q=kl-\tfrac14(s_0l-t_1k)^2.
\]
The transverse quadratic form along \(k=l=0\) has unit
Hessian. Its equivariant factorization produces the associated
node chart \(UV=q\), with the actual auxiliary line character.
At \(q=0\) this compactification carries an additional
vertical fixed surface. Therefore identifying the complete
fixed trace with the horizontal blowup alone has not been
proved. The missing surface and its end must be included
before those local indices can be used as global values.

The residual ruling end is the actual
\(S^2\times S^3\) clutch
\(\tau(h,q_H)=(q_H^{-1}h q_H,q_H)\), with its typed inverse
in the forward period direction. Its mapping torus is the
projective line bundle over \(S^3\times S^1\) of the
quaternion clutch, with \(c_2=\pm1\) according to that
unresolved boundary placement. The charge-\(a\) line gives
the exactly computed pair \((-a^3c_2,-4ac_2)\) and index
\(a(a^2-1)c_2/6\). These computations record the dependence
on the actual boundary orientation; they do not select its
placement in the complete bordism.

The full finite cap arrow is now explicit. In its root
Cartier frame,
\[
 U_{\rm cap}=z^{-ma/2}U_a,\qquad
 B=2(-1)^{b+1}\lambda h U_{\rm cap}^{b},\qquad
 h=1-T_0q_\infty.
\]
For \((a,b)=(2,1)\) it commutes with the actual finite cap
map and retains one forward \(t\)-to-\(s\) ruling arrow.
For \((a,b)=(1,2)\), its exact commutator changes the finite
character by \(m/2\), and the second step closes it. These
intermediate objects and the complete annular factors are
retained in
\path{notes/part3_source_zero_annular_attachment_working.md}.

The active construction cuts the compact suspension outside
the short fixed-normalization graph, removes a fixed compact
Cox-origin neighbourhood, and uses the finite Cox moment
cobordism on that cut. The source translation, its finite-cap
arrow above and the associated-node chart have been written.
On the stable outside portion the newly constructed exact
metric arrow is
\[
 N_t=|\xi|\|t\|^a=R^aN_s,\quad R=\|s\|\|t\|,\qquad
 (\xi,\xi_{\rm dual})\longmapsto
 (R^a\xi,R^{-a}\xi_{\rm dual}).
\]
It preserves \(UV=q\), both node branches and the physical
circle on \(R\ge\varepsilon\). The Tate index change
\(j\mapsto j+1\) absorbs exactly the original \(q\)
factor. At the inner boundary the fixed trace is the actual
\(S^3\), with coordinate \(p=e^{i\phi}(S_0,S_1)\) and
positive spinor frame
\((e^{i(a-1)\phi}z,e^{ia\phi}\overline V)\).
The remaining local map is the surgery joining this framed
fixed boundary to the stable outside end, with its physical
quotient line and outward orientation. In the reduced case
the positive node model is the explicitly identified total
space of \(\mathcal O(-1)\oplus\mathcal O(-2)\) over
\(\mathbb P^2\); its fixed surface and isolated
charge-one point are both retained. Its boundary has not
been substituted for the surgered end. The earlier \(j=-1\)
designation used the SZ chart index. The corresponding
identification with the current labelled fan requires the
actual node-to-Tate character map and its Cartier line;
that index alone supplies no such identification. The
full source representative \(Z=P/S\), including the
winding of \(P\), is retained in the active comparison.

There is also an exact obstruction to replacing the present
end by a radial sphere link. In the charge-one chart,
\[
 R^2=(|S_0|^2+|S_1|^2)(|zV|^2+|T_1|^2).
\]
Thus \(R\) vanishes on the entire fibre over
\([0:0:1]\), not only on its zero section. The verified
local fixed-boundary surgery therefore requires an actual
comparison to this product-radius end. Its closed
projective spinor cap has pair
\((3a^2-3a+1,4)\) and index \(a(a-1)/2\);
the charge-one point has pair \((-1,-4)\) and index zero.
The oriented local surgered end on the reduced branch
has pair \((-8,-8)\) and index \(-1\), as proved in
RT.10 of the real-trace working record and independently
checked there. This remains a local statement. Removing
the capped exceptional disk from the global fixed trace
also removes its intersection class; it cannot be counted
again in the complementary capped trace. The labelled
global attachment, including that cap assignment, remains
under construction.

The alternative anti-holomorphic Tate filling retains its
boundary parameter \(\kappa=4\lambda^2/w\) and replaces
it on the filling disc by
\(4\lambda^2\overline w/R^2\), with a positive radial
collar interpolation. This matches the displayed boundary
parameter and stays in the contraction domain. Its full
source-line-cap square and global orientation have not
been supplied, so it also has not assigned \(e\).

\par\medskip\noindent\textbf{The other remaining return sectors.}\par\nopagebreak[4]
For OS49 the complete actual primitive periods in K65 have
windings \((0,-m,0,2m)\) on the positive \(r\)-loop.
Its two diagonal values and shorter endpoint are proved.
K76 extends the relative-line proof of K72 to the two
bottom-row windings \((L,J)=(-m,2m)\).
The actual boundary connection difference remains
\(-n(du+\alpha\,dv)\); the additional \(L\)-term enters
the complete collar cube. The completed integer comparison,
with its exact same-closing-exponent local calibration, is in
\path{notes/part3_bottom_row_winding_working.md}.
It gives the actual vertical-fan integers
\(-m(n-1)^2\), or \(-k(n-1)^2\) on the shorter loop.
It does not choose a free integer from the earlier residue.
K77 completes the nonvertical construction with its full
period and Cox maps, its integer target-prism solution,
and all intermediate cap and shorter-endpoint arrows.
It also proves the exact image on every cap order and
the retained marked cap-data loop calculation. The
added ordinary-clutch sectors remain on the list above.

The source-loop catalogue, its discriminant and braid maps,
and the endpoint constraints remain in the named-loop and
source-moduli working records. K71 now computes all fans of
the specified selected-cusp loop; K73 and K74 compute all
fans of the specified whole-zero scalar loops and their
retained marked cap/site variations. The remaining source
generators require their own closed source and cap arrows.
For arbitrary unmarked endpoints the endpoint automorphism
must be carried through the lattice, physical circle and
finite-cap data before its return can be combined with
these values. The corresponding sectors have not been
assigned a value by changing the marking silently.
\end{proof}
